\SourceLeaf{212}
\Needspace{12\baselineskip}
\begin{center}
\textit{Formae sub aliis contentae quibus tamen non aequivalent.}
\end{center}

\Needspace{8\baselineskip}
\Article{213}
Supra artt. 157, 158 ostendimus, si forma $f$ determinantis $D$ formam $F$ determinantis $E$ implicet atque in ipsam transeat per substitutionem $\alpha,\beta,\gamma,\delta$, fore $E=(\alpha\delta-\beta\gamma)^2D$; si fuerit $\alpha\delta-\beta\gamma=\pm1$, formam $f$ non modo implicare formam $F$ sed ipsi aequivalentem esse et proin si $f$ ipsam $F$ implicet neque vero eidem aequivaleat, quotientem $\frac{E}{D}$ esse integrum maiorem quam 1. Problema itaque hic solvendum erit, \textit{diiudicare an forma data $f$ determinantis $D$ formam datam $F$ determinantis $Dee$ implicet}, ubi $e$ supponitur esse numerus positivus maior quam 1. Hoc negotium ita absolvemus, ut multitudinem finitam formarum sub $f$ contentarum assignare doceamus quae ita sint comparatae, ut $F$ si sub $f$ contenta est necessario alicui ex illis aequivalere debeat.

I. Ponamus omnes divisores (positivos) numeri $e$ (inclusis etiam 1 et $e$) esse $m,m',m''$ etc., atque $e=mn=m'n'=m''n''$ etc. Designemus brevitatis gratia formam in quam $f$ transit per substitutionem propriam $m.\,0,0,n$ ita $(m;0)$, formam in quam $f$ transit per substitutionem propriam $m,1,0,n$ per $(m;1)$ etc. generaliterque formam in quam $f$ per subst. propriam, $m,k,0,n$ transmutatur per $(m;k)$. Simili modo transeat $f$ per subst. propriam $m',0,0,n'$ in $(m';0)$; per hanc $m',1,0,n'$ in $(m';1)$ etc., per $m'',0,0,n''$ in $(m'';0)$ etc. etc. Omnes hae formae sub $f$ proprie contentae erunt, et cuiusvis determinans $=Dee$. Complexum omnium formarum $(m;0)$, $(m;1)$, $(m;2)\ldotp\ldotp\ldotp\ldotp(m;m-1)$, $(m';0)$, $(m';1)\ldots(m';m'-1)$; $(m'';0)$ etc. quarum multitudo erit $m+m'+m''+\text{etc.}$ et quas omnes inter se diversas fore facile perspicitur, designemus per $\Omega$.

Si \textit{e. g.} forma $f$ est haec $(2,5,7)$ atque $e=5$, $\Omega$ comprehendet sequentes sex formas $(1;0)$, $(5;0)$, $(5;1)$, $(5;2)$, $(5;3)$, $(5;4)$ quae si evolvuntur sunt $(2,25,175)$, $(50,25,7)$, $(50,35,19)$, $(50,45,35)$, $(50,55,55)$, $(50,65,79)$.

II. Iam dico, si forma $F$ determinantis $Dee$ sub $f$ proprie contenta sit, necessario eandem alicui formarum $\Omega$ proprie aequivalentem fore. Ponamus formam $f$ transformari in $F$ per substitutionem propriam $\alpha,\beta,\gamma,\delta$, eritque $\alpha\delta-\beta\gamma=e$. Sit numerorum $\gamma,\delta$ (qui ambo simul 0 esse nequeunt) divisor communis maximus positive acceptus $=n$, atque $\frac{e}{n}=m$, qui manifesto erit in%
\SourceLeaf{213}%
teger. Accipiantur $g,h$ ita ut sit $\gamma g+\delta h=n$, denique sit $k$ residuum minimum positivum numeri $\alpha g+\beta h$ secundum modulum $m$. Tum forma $(m;k)$ quae manifesto erit inter formas $\Omega$, formae $F$ proprie aequivalebit, et quidem in ipsam transformabitur per substitutionem propriam
\[
\frac{\gamma}{n}\cdot\frac{\alpha g+\beta h-k}{m}+h,\quad
\frac{\delta}{n}\cdot\frac{\alpha g+\beta h-k}{m}-g,\quad
\frac{\gamma}{n},\quad\frac{\delta}{n}
\]
Nam \textit{primo} perspicuum est hos quatuor numeros esse integros; \textit{secundo} facile confirmatur substitutionem esse propriam; \textit{tertio} patet, formam in quam $(m;k)$ per substitutionem illam transeat eandem esse in quam $f$\HistoricalNote{*)}{Quippe quae per substitutionem $m,k,0,n$ in $(m;k)$ transit V. art. 159.} transeat per substitutionem
\[
m\left(\frac{\gamma}{n}\cdot\frac{\alpha g+\beta h-k}{m}+h\right)+\frac{k\gamma}{n},\quad
m\left(\frac{\delta}{n}\cdot\frac{\alpha g+\beta h-k}{m}-g\right)+\frac{k\delta}{n},\quad\gamma,\quad\delta
\]
sive quoniam $mn=e=\alpha\delta-\beta\gamma$, adeoque $\beta\gamma+mn=\alpha\delta$, $\alpha\delta-mn=\beta\gamma$, per hanc
\[
\frac{1}{n}(\alpha\gamma g+\alpha\delta h),\quad
\frac{1}{n}(\beta\gamma g+\beta\delta h),\quad\gamma,\quad\delta
\]
sive denique quoniam $\gamma g+\delta h=n$, per hanc $\alpha,\beta,\gamma,\delta$ \textit{i. e.} per hyp., in $F$. Quare $(m;k)$ et $F$ proprie aequivalentes erunt. \textit{Q. E. D.}

Ex his igitur semper diiudicari potest, an forma aliqua data $f$ determinantis $D$ formam $F$ determinantis $Dee$ proprie implicet. Si vero quaeritur an $f$ ipsam $F$ improprie implicet, investigari tantummodo debet an forma ipsi $F$ opposita sub $f$ proprie contenta sit, art. 159.

\Needspace{8\baselineskip}
\Article{214}
\textsc{Problema.} \textit{Propositis duabus formis, $f$ determinantis $D$, et $F$ determinantis $Dee$, quarum prior posteriorem proprie implicat: exhibere omnes transformationes proprias formae $f$ in $F$.}

\textit{Sol.} Designante $\Omega$ eundem formarum complexum ut in art. praec., excerpantur ex hoc complexu omnes formae quibus $F$ proprie aequivalet, quae sint $\Phi,\Phi',\Phi''$ etc. Quaevis harum formarum sequenti modo suppeditabit transformationes proprias formae $f$ in $F$, et quidem aliae alias (\textit{i. e.} singulae diversas), cunctae vero cunctas (\textit{i. e.} nulla transformatio propria formae $f$ in $F$ erit quam non una ex formis $\Phi,\Phi'$ etc. praebeat). Quoniam methodus pro omnibus formis $\Phi,\Phi'$ etc. eadem est, de una tantum loquemur.

\SourceLeaf{214}
Ponamus $\Phi$ esse $(M;K)$, atque $e=MN$ ita ut $f$ in $\Phi$ per substitutionem propriam $M,K,0,N$ transeat. Porro designentur omnes transformationes propriae formae $\Phi$ in $F$ indefinite per $\mathfrak{a},\mathfrak{b},\mathfrak{c},\mathfrak{d}$. Tum manifesto $f$ transibit in $\Phi$ per substitutionem propriam $M\mathfrak{a}+K\mathfrak{c},M\mathfrak{b}+K\mathfrak{d},N\mathfrak{c},N\mathfrak{d}$, et hoc modo ex quavis transformatione propria formae $\Phi$ in $F$ sequetur transformatio propria formae $f$ in $F$. --- Eodem modo tractandae sunt formae reliquae $\Phi',\Phi''$ etc., quarum singulae transformationes propriae in $F$ transformationem propriam formae $f$ in $F$ praebebunt.

Ut appareat, hanc solutionem ex omni parte completam esse, ostendendum erit

I. \textit{Hoc modo omnes transformationes proprias possibiles formae $f$ in $F$ obtineri.} Sit transformatio quaecunque propria formae $f$ in $F$ haec $\alpha,\beta,\gamma,\delta$ atque ut in art. praec. II, $n$ divisor communis maximus numerorum $\gamma,\delta$; numeri $m,g,h,k$ autem eodem modo ut illic determinati. Tunc forma $(m;k)$ erit inter formas $\Phi,\Phi'$ etc., et
\[
\frac{\gamma}{n}\cdot\frac{\alpha g+\beta h-k}{m}+h,\quad
\frac{\delta}{n}\cdot\frac{\alpha g+\beta h-k}{m}-g,\quad
\frac{\gamma}{n},\quad\frac{\delta}{n}
\]
aliqua ex transformationibus propriis huius formae in $F$; ex hac vero per regulam modo traditam obtinetur transformatio $\alpha,\beta,\gamma,\delta$; haec omnia in art. praec. sunt demonstrata.

II. \textit{Omnes transformationes hoc modo prodeuntes inter se diversas esse, seu nullam bis obtineri.} Nullo quidem negotio perspicitur, plures transformationes diversas eiusdem formae $\Phi$ vel $\Phi'$ etc. in $F$ eandem transformationem formae $f$ in $F$ producere non posse; quod vero etiam formae diversae e. g. $\Phi$ et $\Phi'$ eandem transformationem suppeditare nequeant, ita demonstratur. Supponamus, transformationem propriam $\alpha,\beta,\gamma,\delta$ formae $f$ in $F$ obtineri \textit{tum} ex transformatione propria $\mathfrak{a},\mathfrak{b},\mathfrak{c},\mathfrak{d}$ formae $\Phi$ in $F$, \textit{tum} ex transformatione propria $\mathfrak{a}',\mathfrak{b}',\mathfrak{c}',\mathfrak{d}'$ formae $\Phi'$ in $F$. Sit $\Phi=(M;K)$, $\Phi'=(M';K')$, $e=MN=M'N'$. Habebuntur itaque aequationes
\begin{align*}
\alpha&=M\mathfrak{a}+K\mathfrak{c}=M'\mathfrak{a}'+K'\mathfrak{c}' \tag*{. . . . . . . [1]}\\
\beta&=M\mathfrak{b}+K\mathfrak{d}=M'\mathfrak{b}'+K'\mathfrak{d}' \tag*{. . . . . . . [2]}\\
\gamma&=N\mathfrak{c}=N'\mathfrak{c}' \tag*{. . . . . . . . . . . . . [3]}
\\
\SourceLeaf{215}\delta&=N\mathfrak{d}=N'\mathfrak{d}' \tag*{. . . . . . . . . . . . . [4]}\\
\mathfrak{a}\mathfrak{d}-\mathfrak{b}\mathfrak{c}&=\mathfrak{a}'\mathfrak{d}'-\mathfrak{b}'\mathfrak{c}'=1 \tag*{. . . . . . . . . [5]}
\end{align*}
Ex $\mathfrak{a}[4]-\mathfrak{b}[3]$ sequitur adiumento aequ. [5], $N=N'(\mathfrak{a}\mathfrak{d}'-\mathfrak{b}\mathfrak{c}')$, quare $N'$ metietur ipsum $N$; similiter ex $\mathfrak{a}'[4]-\mathfrak{b}'[3]$ fit $N(\mathfrak{a}'\mathfrak{d}-\mathfrak{b}'\mathfrak{c})=N'$, quare $N$ metietur ipsum $N'$, unde, quia tum $N$ tum $N'$ supponuntur esse positivi, erit necessario $N=N'$, et $M=M'$, et hinc ex 3 et 4, $\mathfrak{c}=\mathfrak{c}'$, $\mathfrak{d}=\mathfrak{d}'$. Porro fit ex $\mathfrak{a}[2]-\mathfrak{b}[1]$,
\[
K=M'(\mathfrak{a}\mathfrak{b}'-\mathfrak{b}\mathfrak{a}')+K'(\mathfrak{a}\mathfrak{d}'-\mathfrak{b}\mathfrak{c}')=M(\mathfrak{a}\mathfrak{b}'-\mathfrak{b}\mathfrak{a}')+K'
\]
hinc $K\equiv K'\;(\mathrm{mod.}\,M)$ quod fieri nequit nisi $K=K'$, quia tum $K$ tum $K'$ iacent inter limites 0 et $M-1$. Quamobrem formae $\Phi,\Phi'$ non sunt diversae, contra hyp.

Ceterum patet, si $D$ fuerit negativus vel positivus quadratus, per methodum hanc omnes transformationes proprias formae $f$ in $F$ revera inveniri posse; si vero $D$ positivus non-quadratus; formulae certae generales assignari poterunt, in quibus omnes transformationes propriae (quarum multitudo infinita) contentae erunt.

Denique, si forma $F$ improprie sub forma $f$ contenta est, omnes transformationes impropriae illius in hanc per methodum traditam facile exhiberi poterunt. Scilicet si $\alpha,\beta,\gamma,\delta$ indefinite omnes transformationes proprias formae $f$ in formam quae formae $F$ opposita est, designare supponitur: omnes transf. impropriae formae $f$ in $F$ exhibebuntur per $\alpha,-\beta,\gamma,-\delta$.

\textit{Ex.} Desiderantur omnes transformationes formae $(2,5,7)$ in $(275,0,-1)$, quae sub illa tum proprie tum improprie contenta est. Complexum formarum $\Omega$ pro hoc casu iam in art. praec. tradidimus; examine instituto invenitur, tum $(5;1)$ tum $(5;4)$ formae $(275,0,-1)$ proprie aequivalere. Omnes transformationes propriae formae $(5;1)$ \textit{i. e.} $(50,35,19)$ in $(275,0,-1)$ per theoriam nostram supra explicatam inveniuntur contineri sub formula generali
\[
16t-275u,\quad -t+16u,\quad -15t+275u,\quad t-15u
\]
ubi $t,u$ designant indefinite omnes numeros integros aequationi $tt-275uu=1$
\SourceLeaf{216}
satisfacientes; quare omnes transformationes propriae formae $(2,5,7)$ in $(275,0,-1)$ hinc oriundae contentae erunt sub formula generali
\[
65t-1100u,\quad -4t+65u,\quad -15t+275u,\quad t-15u
\]
Simili modo omnes transformationes propriae formae $(5,4)$ \textit{i. e.} $(50,65,79)$ in $(275,0,-1)$ continentur sub formula generali
\[
14t+275u,\quad t+14u,\quad -15t-275u,\quad -t-15u
\]
adeoque omnes transformationes propriae formae $(2,5,7)$ in $(275,0,-1)$ hinc oriundae sub hac
\[
10t+275u,\quad t+10u,\quad -15t-275u,\quad -t-15u
\]
Hae duae formulae igitur omnes transformationes proprias quaesitas amplectuntur\HistoricalNote{*)}{Concinnius omnes transformationes propriae exhibentur per formulam
\[
10t+55u,\quad t+2u,\quad -15t-55u,\quad -t-3u
\]
denotantibus $t,u$ indefinite omnes integros aequationi $tt-11uu=1$ satisfacientes.}. --- Eodem vero modo invenitur, omnes transformationes improprias formae $(2,5,7)$ in $(275,0,-1)$ sub sequentibus duabus formulis contentas esse:
\[
\text{(I)}\;\ldots\quad 65t-1100u,\quad 4t-65u,\quad -15t+275u,\quad -t+15u
\]
et
\[
\text{(II)}\;\ldots\quad 10t+275u,\quad -t-10u,\quad -15t-275u,\quad t+15u
\]

\Needspace{12\baselineskip}
\begin{center}
\textit{Formae determinantis 0.}
\end{center}

\Needspace{8\baselineskip}
\Article{215}
Hucusque formas determinantis 0 ab omnibus disquisitionibus exclusimus; de his itaque, ut theoria nostra ab omni parte completa evadat, quaedam adhuc sunt adiicienda. Quoniam generaliter demonstratum est, si forma aliqua determinantis $D$ formam determinantis $D'$ implicet, $D'$ esse multiplum ipsius $D$, statim patet, formam cuius determinans $=0$ aliam formam quam cuius determinans etiam sit $=0$ implicare non posse. Quare duo tantummodo problemata solvenda restant, scilicet $1^{\mathrm{o}}$ \textit{propositis duabus formis $f,F$, quarum posterior habet determinantem 0, diiudicare utrum prior posteriorem implicet necne, et in illo casu omnes transformationes illius in hanc exhibere.} $2^{\mathrm{o}}$ \textit{Invenire omnes repraesentationes numeri dati per formam datam determinantis 0.} Problema primum aliam metho%
\SourceLeaf{217}%
dum requirit, quando determinans prioris formae $f$ etiam est 0, aliam quando non est 0. Haec omnia iam exponemus.

I. Ante omnia observamus, quamvis formam $axx+2bxy+cyy$, cuius determinans $bb-ac=0$, ita exhiberi posse $m(gx+hy)^2$, denotantibus $g,h$ numeros inter se primos, $m$ integrum. Sit enim $m$ divisor communis maximus ipsorum $a,c$ eodem signo acceptus quo hi numeri ipsi sunt affecti (hos signa opposita habere non posse facile perspicitur), eruntque $\frac{a}{m},\frac{c}{m}$ integri inter se primi non negativi, productum ex ipsis $=\frac{bb}{mm}$ \textit{i. e.} quadratum, adeoque illi ipsi quadrata (art. 21). Sit $\frac{a}{m}=gg$, $\frac{c}{m}=hh$, eruntque etiam $g,h$ inter se primi, $gghh=\frac{bb}{mm}$, et $gh=\pm\frac{b}{m}$. Hinc patet
\[
m(gx\pm hy)^2\text{ fore }=axx+2bxy+cyy
\]

Iam propositae sint duae formae $f,F$, utraque determinantis 0, et quidem sit
\[
f=m(gx+hy)^2,\quad F=M(GX+HY)^2
\]
ita ut $g$ ad $h$, $G$ ad $H$ sint primi. Tum dico, si forma $f$ implicet formam $F$, $m$ aut ipsi $M$ aequalem esse aut saltem ipsum $M$ metiri et quotientem esse quadratum; et vice versa si $\frac{M}{m}$ sit quadratum integrum, $F$ contentam esse sub $f$. Si enim $f$ per substitutionem
\[
x=\alpha X+\beta Y,\quad y=\gamma X+\delta Y
\]
in $F$ transire supponitur, erit
\[
\frac{M}{m}(GX+HY)^2=((\alpha g+\gamma h)X+(\beta g+\delta h)Y)^2
\]
unde facile sequitur $\frac{M}{m}$ esse quadratum. Ponatur $=ee$, eritque
\[
e(GX+HY)=\pm((\alpha g+\gamma h)X+(\beta g+\delta h)Y),\quad\text{\textit{i. e.}}
\]
\[
\pm eG=\alpha g+\gamma h,\quad\pm eH=\beta g+\delta h
\]
si itaque $\mathfrak{G},\mathfrak{H}$ ita determinantur ut sit $\mathfrak{G}G+\mathfrak{H}H=1$, erit
\[
\pm e=\mathfrak{G}(\alpha g+\gamma h)+\mathfrak{H}(\beta g+\delta h),\quad\text{adeoque \textit{integer.}\quad\textit{Q. E. P.}}
\]
Si vero, vice versa, supponitur, $\frac{M}{m}$ esse quadratum integrum $=ee$, forma $f$ implicabit formam $F$. Scilicet integri $\alpha,\beta,\gamma,\delta$ ita poterunt determinari ut fiat
\SourceLeaf{218}
\[
\alpha g+\gamma h=\pm eG,\quad \beta g+\delta h=\pm eH
\]
Accipiantur enim integri $\mathfrak{g},\mathfrak{h}$ ita ut fiat $\mathfrak{g}g+\mathfrak{h}h=1$, satisfietque aequationibus illis ponendo
\[
\begin{aligned}
\alpha&=\pm eG\mathfrak{g}+hz, &\quad\gamma&=\pm eG\mathfrak{h}-gz\\
\beta&=\pm eH\mathfrak{g}+hz', &\quad\delta&=\pm eH\mathfrak{h}-gz'
\end{aligned}
\]
quicunque valores integri ipsis $z,z'$ tribuantur; quare $F$ contenta erit sub $f$, \textit{Q. E. S.} Simul haud difficulter intelligitur, has formulas omnes valores quos $\alpha,\beta,\gamma,\delta$ nancisci possunt, \textit{i. e.} omnes transformationes formae $f$ in $F$ exhibere, si modo $z,z'$ indefinite omnes numeros integros exhibere supponantur.

II. Propositis duabus formis $f=axx+2bxy+cyy$ cuius determinans non $=0$, et $F=M(GX+HY)^2$ cuius determinans $=0$ (designantibus ut ante $G,H$ numeros inter se primos), dico \textit{primo}, si $f$ implicet ipsam $F$, numerum $M$ per formam $f$ repraesentari posse; \textit{secundo}, si $M$ per $f$ repraesentari possit, $F$ sub $f$ contentam esse; \textit{tertio}, si in hoc casu omnes repraesentationes numeri $M$ per formam $f$ indefinite exhibeantur ita $x=\xi,y=\upsilon$, omnes transformationes formae $f$ in $F$ exhiberi ita $G\xi,H\zeta,G\upsilon,H\upsilon$. Quae omnia sequenti modo demonstramus.

$1^{\mathrm{o}}$ Ponamus $f$ transire in $F$ per substitutionem $\alpha,\beta,\gamma,\delta$ accipianturque numeri $\mathfrak{G},\mathfrak{H}$ ita ut sit $\mathfrak{G}G+\mathfrak{H}H=1$. Tunc manifestum est, si ponatur $x=\alpha\mathfrak{G}+\beta\mathfrak{H}$, $y=\gamma\mathfrak{G}+\delta\mathfrak{H}$, valorem formae $f$ fieri $M$, adeoque $M$ repraesentabilem esse per formam $f$.

$2^{\mathrm{o}}$ Si supponitur esse $a\xi\xi+2b\xi\upsilon+c\upsilon\upsilon=M$, manifestum est, per substitutionem $G\xi,H\xi,G\upsilon,H\upsilon$, formam $f$ transire in $F$. Quod vero

$3^{\mathrm{o}}$ in hoc casu substitutio $G\xi,H\zeta,G\upsilon,H\upsilon$ \textit{omnes} transformationes formae $f$ in $F$ exhibeat, si $\xi,\upsilon$ supponantur exhibere omnes valores ipsorum $x,y$, qui faciunt $f=M$, ita perspicitur. Sit $\alpha,\beta,\gamma,\delta$ transformatio quaecunque formae $f$ in $F$, et ut ante $\mathfrak{G}G+\mathfrak{H}H=1$. Tum inter valores ipsorum $x,y$ erunt etiam hi
\[
x=\alpha\mathfrak{G}+\beta\mathfrak{H},\quad y=\gamma\mathfrak{G}+\delta\mathfrak{H}
\]
ex quibus obtinetur substitutio
\SourceLeaf{219}
\[
G(\alpha\mathfrak{G}+\beta\mathfrak{H}),\quad H(\alpha\mathfrak{G}+\beta\mathfrak{H}),\quad G(\gamma\mathfrak{G}+\delta\mathfrak{H}),\quad H(\gamma\mathfrak{G}+\delta\mathfrak{H})
\]
sive
\[
\begin{aligned}
&\alpha+\mathfrak{H}(\beta G-\alpha H), &\quad&\beta+\mathfrak{G}(\alpha H-\beta G)\\
&\gamma+\mathfrak{H}(\delta G-\gamma H), &&\delta+\mathfrak{G}(\gamma H-\delta G)
\end{aligned}
\]
Sed quoniam
\[
\begin{aligned}a(\alpha X+\beta Y)^2+2b(\alpha X+\beta Y)(\gamma X+\delta Y)&{}\\
{}+c(\gamma X+\delta Y)^2&=M(GX+HY)^2\end{aligned}
\]
erit
\[
\begin{aligned}
a(\alpha\delta-\beta\gamma)^2&=M(\delta G-\gamma H)^2\\
c(\beta\gamma-\alpha\delta)^2&=M(\beta G-\alpha H)^2
\end{aligned}
\]
adeoque (quum determinans formae $f$ per $(\alpha\delta-\beta\gamma)^2$ multiplicatus aequalis sit determinanti formae $F$ \textit{i. e.} $=0$, adeoque etiam $\alpha\delta-\beta\gamma=0$),
\[
\delta G-\gamma H=0,\quad \beta G-\alpha H=0
\]
Hinc substitutio illa transit in hanc $\alpha,\beta,\gamma,\delta$, unde patet, formulam traditam \textit{omnes} transformationes formae $f$ in $F$ suppeditare.

III. Superest ut omnes repraesentationes numeri dati per formam datam determinantis 0 exhibere doceamus. Sit forma haec $m(gx+hy)^2$, patetque statim, numerum illum per $m$ divisibilem, et quotientem quadratum esse debere. Si itaque numerus propositus statuitur $=mee$, perspicuum est, pro quibus valoribus ipsorum $x,y$ fiat $m(gx+hy)^2=mee$, pro iisdem fieri $gx+hy$ aut $=+e$, aut $=-e$. Quare omnes repraesentationes habebuntur, si omnes solutiones aequationum linearium $gx+hy=e$, $gx+hy=-e$ in integris. sunt inventae. Has vero solubiles esse constat (siquidem $g,h$ sunt inter se primi ut supponitur). Scilicet si $\mathfrak{g},\mathfrak{h}$ ita determinantur ut sit $\mathfrak{g}g+\mathfrak{h}h=1$, aequationi priori satisfiet ponendo $x=\mathfrak{g}e+hz$, $y=\mathfrak{h}e-gz$; posteriori vero faciendo $x=-\mathfrak{g}e+hz$, $y=-\mathfrak{h}e-gz$, denotante $z$ integrum quemcunque. Simul vero formulae hae \textit{omnes} valores integros ipsorum $x,y$ exhibent, si $z$ indefinite numerum quemvis integrum designare supponitur.

\SourceLeaf{220}
\Needspace{12\baselineskip}
\begin{center}
\textit{Solutio generalis omnium aequationum indeterminatarum secundi gradus duas incognitas implicantium per numeros integros.}
\end{center}

His disquisitionibus coronidis loco apponimus

\Needspace{8\baselineskip}
\Article{216}
\textsc{Problema.} \textit{Invenire omnes solutiones aequationis generalis}\HistoricalNote{*)}{Si aequatio proponeretur in qua coefficiens secundus, quartus vel quintus non esset par, multiplicata per 2 eam formam reciperet quam hic supponimus.} \textit{indeterminatae secundi gradus duas incognitas implicantis}
\[
axx+2bxy+cyy+2dx+2ey+f=0
\]
\textit{(ubi $a,b,c$ etc. sunt integri quicunque dati) per numeros integros.}

\textit{Sol.} Introducamus loco incognitarum $x,y$ alias
\[
p=(bb-ac)x+be-cd\quad\text{et}\quad q=(bb-ac)y+bd-ae
\]
qui manifesto semper erunt integri, quando $x,y$ sunt integri. Quo facto habebitur aequatio
\[
app+2bpq+cqq+f(bb-ac)^2+(bb-ac)(aee-2bde+cdd)=0
\]
sive posito brevitatis gratia numero
\[
f(bb-ac)^2+(bb-ac)(aee-2bde+cdd)=-M
\]
haec
\[
app+2bpq+cqq=M
\]
Iam omnes solutiones huius aequationis, \textit{i. e.} omnes repraesentationes numeri $M$ per formam $(a,b,c)$ in praecedentibus invenire docuimus. Si vero ex singulis valoribus ipsorum $p,q$, valores respondentes ipsorum $x,y$ adiumento aequationum
\[
x=\frac{p+cd-be}{bb-ac},\quad y=\frac{q+ae-bd}{bb-ac}
\]
determinantur, facile perspicitur, omnes hos valores aequationi propositae satisfacere, et nullos valores integros ipsorum $x,y$ dari qui hoc modo non obtineantur. Si itaque ex omnibus valoribus ipsorum $x,y$ sic prodeuntibus valores fractos eiicimus, omnes solutiones quaesitae remanebunt.

Circa hanc solutionem sequentia sunt observanda.

\SourceLeaf{221}
$1^{\mathrm{o}}$ Si aut $M$ per formam $(a,b,c)$ repraesentari non potest, aut ex nulla repraesentatione valores \textit{integri} ipsorum $x,y$ sequuntur: aequatio in integris nullo modo solvi poterit.

$2^{\mathrm{o}}$ Quando determinans formae $(a,b,c)$, \textit{i. e.} numerus $bb-ac$ est negativus, vel positivus quadratus simulque $M$ non $=0$: multitudo repraesentationum numeri $M$ per formam $(a,b,c)$ erit finita, et proin etiam multitudo omnium solutionum aequationis propositae (si quae omnino dantur) finita erit.

$3^{\mathrm{o}}$ Quando $bb-ac$ est positivus non-quadratus, vel quadratus et simul $M=0$: numerus $M$, si ullo modo, \textit{infinitis modis diversis} per formam $(a,b,c)$ repraesentari poterit; sed quoniam impossibile est, has repraesentationes omnes \textit{ipsas} invenire et tentare utrum valores integros ipsorum $x,y$ praebeant an fractos, necessarium est regulam tradere, per quam, quando forte nulla omnino repraesentatio valores integros ipsorum $x,y$ praebere potest, de hac re \textit{certi} fieri possimus (nam quotcunque repraesentationes in hoc casu \textit{tentatae} fuerint, absque tali regula ad certitudinem nunquam perveniremus); quando vero aliae repraesentationes dant valores integros ipsorum $x,y$, aliae fractos: docendum erit quomodo hae ab illis a priori generaliter dignosci possint.

$4^{\mathrm{o}}$ Quando $bb-ac=0$: valores ipsorum $x,y$ per formulas praecedentes omnino non possunt determinari; quare pro hoc casu \textit{methodus peculiaris} investigari debebit.

\Needspace{8\baselineskip}
\Article{217}
Pro eo casu, ubi $bb-ac$ est numerus positivus non-quadratus, supra docuimus, omnes repraesentationes numeri $M$ per formam $app+2bpq+cqq$ (si quae omnino dentur) exhiberi posse, per unam vel per plures formulas tales
\[
p=\tfrac{1}{m}(\mathfrak{A}t+\mathfrak{B}u),\quad q=\tfrac{1}{m}(\mathfrak{C}t+\mathfrak{D}u)
\]
denotantibus $\mathfrak{A},\mathfrak{B},\mathfrak{C},\mathfrak{D}$ numeros integros datos, $m$ divisorem communem maximum numerorum $a,2b,c$; denique $t,u$ indefinite omnes numeros integros aequationi $tt-(bb-ac)uu=mm$ satisfacientes. Quoniam omnes valores ipsorum $t,u$ tum positive tum negative accipi possunt: pro singulis illarum formarum \textit{quaternas} alias substituere poterimus,
\[
\begin{aligned}
p&=\tfrac{1}{m}(\mathfrak{A}t+\mathfrak{B}u), &\quad q&=\tfrac{1}{m}(\mathfrak{C}t+\mathfrak{D}u)\\
p&=\tfrac{1}{m}(\mathfrak{A}t-\mathfrak{B}u), &q&=\tfrac{1}{m}(\mathfrak{C}t-\mathfrak{D}u)
\end{aligned}
\]
